\section{Meta-FBOHN v3: Learned Orbit Mixing}
\label{sec:meta_fbohn_v3}

\subsection{Idea}
Variant v3 replaced feature scaling with nonlinear orbit mixing:
\[
Z_j=\gamma_j\tanh\!\left(\frac{1}{\sqrt{30}}\sum_i W_{ji}F_i+b_j\right),
\]
and the classifier received the representation
\[
\Phi_{v3}(x)=[\text{FBOHN}(x),Z_1(x),\ldots,Z_k(x)].
\]

\subsection{Results}
Meta-FBOHN v3 yielded the first small but systematic improvement over FBOHN. Typical values from logs:
\begin{center}
\begin{tabular}{lrr}
\toprule
Model & Number of features & Typical accuracy range\\
\midrule
FBOHN & 30 & 0.61--0.63\\
Meta-FBOHN v3 & 42 & 0.63--0.64\\
FBOHN-poly-control & 84 & 0.63--0.64\\
\bottomrule
\end{tabular}
\end{center}

\subsection{Analysis}
The result was interesting because 12~learned mixes provided a similar improvement to a much larger manual polynomial control. However, v3 features were harder to interpret because they were dense linear combinations passed through a $\tanh$ function.

\subsection{Conclusion}
Meta-FBOHN should not only mix features but discover interpretable symbolic operators.
